\begin{figure*}[htbp]
	\centering
	\includegraphics[width=0.95\textwidth]{figures/overview.pdf}
	\caption{Overview architecture of dataset creation and AutoBool training.}
	\label{fig:overview}
\end{figure*}

\section{Related Work}
\label{sec:related_work}
\textbf{Traditional Approaches to Boolean Query Formulation.}
Previous work on automated Boolean query formulation followed two main paradigms: the \emph{objective method} and the \emph{conceptual method}~\cite{scells2020conceptual, scells2020objective}. The objective method aims to replicate known included studies by expanding on initial seed sets using keyword co-occurrence, term frequency analysis, or relevance feedback~\cite{hausner2012routine}. These approaches prioritize recall but sacrifice interpretability and precision. The conceptual method involves manually identifying key aspects (population, intervention, outcome) and converting them into Boolean expressions. While more interpretable and aligned with expert workflows, they require substantial manual effort and domain expertise~\cite{clark2013systematic}.

\textbf{LLM-Based Query Generation.}
LLMs have recently been explored for automating Boolean query generation from topic descriptions. While they can produce syntactically valid queries, they often require multiple attempts, and their recall remains well below expert-authored queries~\cite{wang2025reassessing}. These limitations highlight the need for trainable methods beyond static prompting. Supervised fine-tuning is not well-suited here, as there is no single high-quality ground-truth query per topic: expert-authored queries are manually refined, inconsistent, and often suboptimal.

\textbf{Reinforcement Learning for LLM Optimization.}
RL offers a flexible way to optimize models based on task-level feedback without requiring gold-standard targets. It has emerged as a promising approach for optimizing LLMs on non-differentiable objectives such as factual accuracy, preference alignment, and retrieval quality~\cite{zhang2020deep, zhuang2025rank, nguyen2024reinforcement}. Notably, Group Relative Policy Optimization (\textbf{GRPO}) has shown strong performance in LLM fine-tuning scenarios, including DeepSeek-R1~\cite{deepseekai2025deepseekv3technicalreport, guo2025deepseek}.
